\begin{abstract}
Large language models (LLMs) are increasingly explored for cyber-physical systems (CPS) security, yet most existing efforts either analyse static datasets, use LLMs only for offline advisory tasks, or evaluate model behaviour without a realistic closed-loop execution environment. This makes it difficult to study whether LLM-driven attack and defence strategies remain effective when exposed to real controller timing, process dynamics, and safety constraints.

This paper presents \textbf{CPSForge}, a closed-loop LLM red-team/blue-team testbed for cyber-physical systems. CPSForge is built around a real hardware-in-the-loop environment consisting of a Siemens S7 PLC programmed in TIA Portal~v17 and connected to Factory~I/O as the physical-process emulator. The framework supports two operational modes: a \emph{batch mode} in which scripted, random, or LLM-based attackers execute structured attack actions through an orchestrated pipeline, and a \emph{live agent mode} in which an autonomous LLM attacker agent and an LLM defender agent operate concurrently as independent processes, with a coordinator mediating all PLC writes through a runtime safety shield. In both modes, every proposed write is validated against configurable safety rules before reaching the PLC, and all execution is logged for reproducible analysis.

We evaluate CPSForge across three Factory~I/O scenes of increasing complexity using four attacker variants---scripted, random, LLM batch, and LLM agent---paired with threshold-based, invariant-based, and LLM-based defenders. Across 55~evaluation runs totalling 4{,}465~live PLC steps, results demonstrate that all attacker variants produce process-valid actions with 100\% execution success, the safety shield blocks only out-of-range random perturbations while permitting all well-formed attacks, and threshold/invariant detectors achieve perfect recall (1.0) across batch-mode configurations. Agent-mode experiments reveal that small local LLMs (3.8B parameters) are insufficient for reliable structured CPS agent operation, establishing a practical lower bound on model capacity for this task. To our knowledge, CPSForge is the first framework to combine LLM-driven attack generation, an autonomous LLM defender, a configurable runtime safety shield, and closed-loop adaptation, all executing against a real PLC in a hardware-in-the-loop setting.
\end{abstract}